%% =========================================================
%%  运动引导双流融合视觉-语言-动作模型
%%  用于长时序机器人操作
%%
%%  编译方式：xelatex main.tex（需安装 ctex 宏包）
%%  图片占位符说明：凡注有 [需图片] 或 [需实验数据] 的
%%  位置，均需在实验完成后替换。
%% =========================================================
\documentclass[10pt,twocolumn]{article}

%% ----- 中文支持 -----
\usepackage[UTF8, zihao=-4]{ctex}

%% ----- 版面 -----
\usepackage[a4paper,
  top=25mm, bottom=25mm,
  left=18mm, right=18mm,
  columnsep=6mm]{geometry}

%% ----- 数学 -----
\usepackage{amsmath,amssymb,mathtools}
\usepackage{bm}

%% ----- 图表 -----
\usepackage{graphicx}
\usepackage{subcaption}
\usepackage{booktabs}
\usepackage{multirow}
\usepackage{float}

%% ----- 颜色与超链接 -----
\usepackage[dvipsnames]{xcolor}
\usepackage[colorlinks=true, linkcolor=NavyBlue,
            citecolor=ForestGreen, urlcolor=NavyBlue]{hyperref}

%% ----- 算法伪代码 -----
\usepackage{algorithm}
\usepackage{algorithmic}

%% ----- 其他 -----
\usepackage{enumitem}
\usepackage{cite}
\usepackage{caption}

\captionsetup{font=small, labelfont=bf}
\setlength{\abovecaptionskip}{4pt}
\setlength{\belowcaptionskip}{0pt}

%% ----- 自定义命令 -----
\newcommand{\eg}{\textit{e.g.},\xspace}
\newcommand{\ie}{\textit{i.e.},\xspace}
\newcommand{\etal}{\textit{et al.}\xspace}
\newcommand{\todo}[1]{\textcolor{red}{\textbf{[TODO: #1]}}}
\newcommand{\needfig}[1]{\begin{center}\fbox{\parbox{0.9\linewidth}{\centering\textcolor{BrickRed}{\textbf{[需图片]} #1}}}\end{center}}
\newcommand{\needtab}[1]{\begin{center}\fbox{\parbox{0.9\linewidth}{\centering\textcolor{NavyBlue}{\textbf{[需实验数据]} #1}}}\end{center}}

%% =========================================================
\begin{document}

%% ----- 标题信息 -----
\title{\textbf{运动引导双流融合的视觉-语言-动作模型\\用于长时序机器人操作}\\[4pt]
  \large Motion-Guided Dual-Stream Fusion for Vision-Language-Action Models\\
  in Long-Horizon Robotic Manipulation}

\author{
  匿名作者\\
  匿名单位\\
  \texttt{anonymous@example.com}
}

\date{}
\maketitle
\thispagestyle{empty}

%% =========================================================
\begin{abstract}
视觉-语言-动作（VLA）模型近年来在通用机器人操作领域取得了显著进展，但现有方法在处理长时序操作任务时仍面临两大核心挑战：（1）多视角图像之间的运动信息未被充分利用，跨视角注意力"盲目"地平等对待所有图像区域；（2）高维动作序列的直接预测导致策略方差偏大，降低了长时序任务的完成率。为此，本文提出\textbf{SimVLA-MGA}（Motion-Guided Attention），在 SimVLA~\cite{simvla2026} 框架之上引入两项核心创新：

\textbf{（i）运动引导跨视角注意力机制}：通过帧差分图驱动的轻量级运动感知卷积网络（MotionCNN），自动生成腕部相机的逐 patch 运动激活图，并将其作为注意力偏置注入双流跨注意力融合模块，使静态视角特征在融合时自动聚焦于运动活跃区域，无需光流标签；

\textbf{（ii）动作变分自编码器（ActionVAE）隐式流匹配策略}：将动作序列编码至低维隐变量空间，在隐空间执行流匹配去噪，再解码为完整动作块，显著降低策略方差。

在 VLABench~\cite{vlabench2025} 基准的 100 类任务上，SimVLA-MGA 相比 SimVLA 基线在整体成功率上提升 \todo{X.X\%}，在长时序复合任务上提升 \todo{X.X\%}，同时仅增加约 3M 参数（相对基线参数量增量 < 0.7\%）。
\end{abstract}

%% =========================================================
\section{引言}
\label{sec:intro}

机器人操作的核心挑战之一是如何将自然语言指令与复杂的多步骤物理操作联系起来。视觉-语言-动作（VLA）模型~\cite{openvla,rt2,pi0,simvla2026} 通过大规模视觉-语言预训练提供了一种有效的解决思路：将预训练视觉语言模型（VLM）作为感知骨干，接入轻量动作头来预测机器人控制指令。

然而，现有 VLA 方法在面对 VLABench~\cite{vlabench2025} 所定义的\textbf{长时序复合任务}时（平均轨迹长度超过 500 步，需要抓取-移动-放置等多阶段协作）仍存在明显不足：

\begin{enumerate}[leftmargin=*, label=\textbf{C\arabic*.}]
  \item \textbf{跨视角融合缺乏运动感知。} 多相机系统中，腕部（wrist）相机直接观测末端执行器的运动过程，包含最丰富的接触与位姿变化信息。但现有双流融合方法（如 DualStreamFusion~\cite{simvla2026}）的跨注意力机制对所有图像区域一视同仁，无法将静态视角的注意力引导至腕部图像中正在发生运动的关键区域。

  \item \textbf{高维动作序列直接建模方差偏大。} 流匹配（Flow Matching）直接在原始动作空间操作，动作序列维度高（$T \times D_a$，$T=10$，$D_a=7$），复合任务中累积误差显著。
\end{enumerate}

\needfig{动机图：左侧展示跨视角注意力"盲目"与"运动引导"的对比热力图；右侧展示在 VLABench 长时序任务上不同方法的成功率曲线。}

本文提出 \textbf{SimVLA-MGA}，在 SimVLA 基线之上以\textbf{最小参数增量}（< 0.7\%）解决上述两个问题：

\textbf{运动引导跨视角注意力（Motion-Guided Cross-Attention, MGCA）}：受 DeltaCNN~\cite{deltacnn2022} 和 MotionDeltaCNN~\cite{motiondeltacnn2023} 的启发，我们提出利用帧差分图提取运动语义。与 DeltaCNN 旨在\textit{省略不变区域的计算}不同，本文的目标是\textit{生成运动激活图来引导注意力}——两者出发点完全不同。轻量 MotionCNN（约 3M 参数，结构化通道剪枝为标准 CNN 的 1/4）将帧差分图映射为每个 patch 的运动强度得分，以偏置形式注入跨注意力 logits；可学习缩放因子从零初始化，训练初期完全退化为原始注意力，学习稳定后自适应调节信号强度。

\textbf{动作 VAE 隐式流匹配（ActionVAE）}：将动作块 $[B, T, D_a]$ 编码至紧凑隐变量 $\mathbf{z} \in \mathbb{R}^{d_z}$（$d_z=32$），在隐空间执行流匹配去噪，再由解码器恢复完整动作序列。隐空间维度比原始动作空间低约 $T \times D_a / d_z \approx 2.2$ 倍，有效降低策略方差。

本文的主要贡献总结如下：
\begin{itemize}[leftmargin=*]
  \item 提出 MotionCNN + 运动引导跨注意力机制，首次将帧差分运动语义引入多视角 VLA 的跨视角特征融合；
  \item 提出 ActionVAE 隐式流匹配策略，在低维隐空间执行扩散/流匹配，降低长时序任务的累积误差；
  \item 在 VLABench 100 类任务上进行全面评测，验证两项创新的有效性及其与双流融合基线的协同增益；
  \item 开源完整代码与模型权重，为社区提供可复现的强基线。
\end{itemize}

%% =========================================================
\section{相关工作}
\label{sec:related}

\subsection{视觉-语言-动作模型}

VLA 模型通过在机器人数据上微调预训练 VLM 来生成动作。RT-2~\cite{rt2} 首次将 VLM 用于端到端策略学习，但推理速度慢。OpenVLA~\cite{openvla} 提出开源 VLA 框架，以离散 token 预测动作。$\pi_0$（pi0）~\cite{pi0} 引入流匹配动作头，生成连续动作块。SimVLA~\cite{simvla2026} 提出以 SmolVLM-500M 为骨干的极简基线，在 LIBERO 上以 0.5B 参数超越多项 $\geq$4B 参数的方法，并严格分析了训练动态对性能的影响。本文在 SimVLA 框架上扩展，针对多视角长时序任务引入运动感知机制。

\subsection{多视角机器人策略}

多相机观测为机器人策略提供了更完整的场景信息。RoboFlamingo~\cite{roboflamingo} 通过视觉特征拼接处理多图像输入。DualStreamFusion 将视角分为静态流（全局场景）和动态流（末端执行器），以跨注意力方式融合，但未引入运动感知。本文在双流框架基础上引入运动激活图，使跨注意力在语义层面感知运动。

\subsection{基于帧差分的运动感知}

利用相邻帧差分提取运动信息是视频理解的经典手段。DeltaCNN~\cite{deltacnn2022}（CVPR 2022）利用帧差分跳过不变区域的 CNN 计算，以节省算力。MotionDeltaCNN~\cite{motiondeltacnn2023}（ICCV 2023）将该思路扩展至移动相机场景（与腕部相机场景高度相关）。与上述工作不同，本文不追求计算稀疏性，而是将帧差分转化为\textbf{运动语义引导信号}，注入 Transformer 的注意力机制，服务于不同的目标。

\subsection{动作序列生成}

扩散策略（Diffusion Policy）~\cite{diffpolicy} 将动作预测建模为去噪过程，流匹配~\cite{flowmatching}（Flow Matching）以更直接的轨迹插值替代扩散过程。RoLD~\cite{rold} 提出在隐空间执行扩散，编码动作块至低维表示再解噪。本文的 ActionVAE 沿用该思路，将其与流匹配结合，并整合至 VLA 框架。

\subsection{VLABench 基准}

VLABench~\cite{vlabench2025}（ICCV 2025）提供 100 类语言条件机器人操作任务，包含原始任务（60 类，平均时长约 120 步）与复合任务（40 类，平均时长超 500 步），评测指标包括进度分（PS）、成功率（SR）、技能召回率（SKR）和参数召回率（PR）。现有最强 VLA 方法在复合任务上成功率仍较低，验证了提升长时序推理能力的必要性。

%% =========================================================
\section{方法}
\label{sec:method}

\subsection{整体框架}
\label{sec:method:overview}

SimVLA-MGA 的整体架构如图~\ref{fig:arch} 所示。给定当前时刻 $t$ 的多视角图像 $\{I_t^v\}_{v=1}^{V}$（$V=4$：front、wrist、image\_0、image\_1）、前一时刻的腕部图像 $I_{t-1}^{\text{wrist}}$、本体感知状态 $\mathbf{q}_t \in \mathbb{R}^7$（$\text{xyz} + \text{euler} + \text{gripper}$）以及自然语言指令 $\ell$，模型输出未来 $T$ 步的动作序列 $\hat{\mathbf{a}}_{t:t+T} \in \mathbb{R}^{T \times 7}$。

\needfig{整体架构图（图1）：5 个 Zone 横向布局，标注所有模块及张量维度。参考 simvla\_architecture.xml。}

整体前向传播分为四个阶段：
\begin{enumerate}[leftmargin=*, label=\textbf{S\arabic*.}]
  \item \textbf{视觉语言编码}：SmolVLM 将所有视角图像与语言指令编码为统一的视觉语言特征序列；
  \item \textbf{运动激活图生成}：MotionCNN 由帧差分图 $\Delta_t$ 生成腕部运动激活图 $\mathbf{M}_t$；
  \item \textbf{运动引导双流融合}：CrossAttentionFusion 以 $\mathbf{M}_t$ 为注意力偏置融合静态流与动态流特征；
  \item \textbf{动作生成}：Action Transformer 结合 ActionVAE 在隐空间执行流匹配，解码输出动作序列。
\end{enumerate}

\subsection{视觉语言骨干网络}
\label{sec:method:vlm}

本文沿用 SimVLA~\cite{simvla2026} 的骨干设计，采用 \textbf{SmolVLM-500M-Instruct}（基于 Idefics3 架构）作为视觉语言编码器。与 SimVLA 不同，VLABench 场景使用 $V=4$ 个视角，我们将所有视角图像按固定顺序拼接后一同送入 VLM。

\paragraph{图像编码}
将输入图像统一缩放至 $384 \times 384$，经 SigLIP 视觉编码器和像素混洗（pixel shuffle）连接层后，每张图像产生 $P=576$ 个视觉 token，维度 $D_{\text{vlm}}=576$。四个视角共产生 $4P$ 个视觉 token。

\paragraph{语言编码}
语言指令经 SmolVLM 分词器编码后，与视觉 token 拼接，送入 Idefics3 语言模型，输出融合后的视觉语言特征 $\mathbf{F} \in \mathbb{R}^{B \times (4P + L) \times D_{\text{vlm}}}$（$L$ 为语言 token 数量）。实践中取最后的语言位置输出作为 VLA 动作头的输入。

\subsection{双流多视角融合}
\label{sec:method:dualstream}

\paragraph{流划分}
受腕部相机直接观测末端执行器运动的特殊性启发，我们将四个视角划分为：
\begin{itemize}[leftmargin=*]
  \item \textbf{静态流（Static Stream）}：front、image\_0、image\_1 的 VLM 特征，记为 $\mathbf{F}_s \in \mathbb{R}^{B \times 3P \times D}$，描述全局场景布局；
  \item \textbf{动态流（Dynamic Stream）}：wrist 的 VLM 特征，记为 $\mathbf{F}_w \in \mathbb{R}^{B \times P \times D}$，描述末端执行器局部状态。
\end{itemize}

\paragraph{跨注意力融合}
以静态流特征为 Query，动态流特征为 Key/Value，执行多头跨注意力：
\begin{equation}
  \mathbf{F}_{\text{fused}} = \text{softmax}\!\left(\frac{\mathbf{Q}\mathbf{K}^\top}{\sqrt{d_k}} + \mathbf{B}_{\text{motion}}\right)\mathbf{V} + \mathbf{F}_s
  \label{eq:crossattn}
\end{equation}
其中 $\mathbf{Q} = \mathbf{F}_s \mathbf{W}_Q$，$\mathbf{K} = \mathbf{F}_w \mathbf{W}_K$，$\mathbf{V} = \mathbf{F}_w \mathbf{W}_V$。$\mathbf{B}_{\text{motion}}$ 为运动引导注意力偏置（见第~\ref{sec:method:mgca} 节）。残差连接确保当偏置为零时退化为原始跨注意力。

\subsection{运动引导跨视角注意力}
\label{sec:method:mgca}

\subsubsection{帧差分运动图}

在连续两步之间，腕部相机图像的像素差分近似捕捉末端执行器及接触物体的运动：
\begin{equation}
  \Delta_t = I_t^{\text{wrist}} - I_{t-1}^{\text{wrist}} \in \mathbb{R}^{B \times 3 \times H \times W}
  \label{eq:delta}
\end{equation}
当 $t=0$（episode 第一步）时，令 $I_{-1}^{\text{wrist}} = I_0^{\text{wrist}}$，使差分图全零，注意力偏置趋于均匀，不影响正常注意力行为。

\subsubsection{MotionCNN 架构}

MotionCNN 将帧差分图 $\Delta_t$ 映射为每个 wrist patch 的运动激活得分 $\mathbf{M} \in [0,1]^{B \times P}$，具体结构如下：

\begin{equation}
  \mathbf{M} = \sigma\!\bigl(\mathbf{W}_o \cdot \text{Flatten}\bigl(\text{AvgPool}_{g\times g}(\phi(\Delta_t))\bigr)\bigr)
  \label{eq:motioncnn}
\end{equation}

其中 $\phi$ 为三层卷积编码器（通道数 $3 \to 16 \to 32 \to 32$，步长 2，结构化剪枝为标准 CNN 通道数的 1/4），$\text{AvgPool}_{g \times g}$ 将特征图池化至 $g \times g = \sqrt{P} \times \sqrt{P}$（与 VLM patch 网格对齐），$\mathbf{W}_o$ 为线性投影，$\sigma$ 为 Sigmoid 激活。

\paragraph{结构化剪枝}
通道数减少为标准值的 1/4，参数量约 3M（< 100MB 显存增量），在保证运动语义提取能力的同时不影响训练吞吐量。

\paragraph{零初始化输出层}
$\mathbf{W}_o$ 和偏置均初始化为零，使得训练开始时 $\mathbf{M} \approx 0.5$（均匀激活），模型先以原始跨注意力为起点学习基础策略，再逐步利用运动信号。

\subsubsection{注意力偏置注入}

将运动激活图作为 Key 维度的注意力偏置注入式~\eqref{eq:crossattn}：
\begin{equation}
  \mathbf{B}_{\text{motion}} = \alpha \cdot \mathbf{M}\,[\text{unsqueeze to } B \times 1 \times 1 \times P]
  \label{eq:bias}
\end{equation}
其中 $\alpha$ 为\textbf{可学习标量}，初始值为 0，广播到所有注意力头和 Query 位置。偏置 $\mathbf{B}_{\text{motion}}$ 在 softmax 之前叠加，使 $\mathbf{M}$ 中得分高的 patch（运动活跃区域）获得更大的注意力权重，引导静态视角特征聚焦于"腕部正在运动的地方"。

图~\ref{fig:motion_vis} 给出运动激活图的定性可视化。

\needfig{图2：运动激活图可视化。左列：腕部当前帧与前一帧；中列：帧差分图 $\Delta_t$；右列：MotionCNN 输出的运动激活图（热力图叠加于原图）。选取抓取、移动、放置三个典型阶段各一例。}

\subsection{隐式扩散动作策略（Latent Diffusion Policy）}
\label{sec:method:ldp}

\subsubsection{动机}

标准流匹配直接在原始动作空间 $\mathbb{R}^{T \times D_a}$（$T=10$，$D_a=7$，共 70 维）上建模速度场。高维动作空间存在两个问题：（1）速度场曲面复杂，网络需要更大容量；（2）在长时序任务中，策略跨越抓取→移动→放置等多阶段，动作块分布方差大、多峰明显，直接回归困难。

受 RoLD~\cite{rold} 启发，我们提出 \textbf{Latent Diffusion Policy（LDP）}：用变分自编码器（ActionVAE）将动作块压缩至 $d_z = 32$ 维紧凑隐变量 $\mathbf{z}$，在低维隐空间执行流匹配（LatentFlowNet），再由条件解码器还原完整动作序列。压缩比为 $70/32 \approx 2.2$，隐空间流形更规则，速度场更易学习。

与 RoLD 的关键区别：（1）RoLD 使用离散 VQVAE，我们采用\textbf{连续 VAE}，梯度直通、与流匹配天然兼容；（2）RoLD 的解码器无语言/视觉条件，我们的解码器\textbf{跨注意力到 VLM 融合特征} $\mathbf{F}_{\text{fused}}$，使动作还原可利用语言语义。

\subsubsection{ActionVAE：编码器}

编码器 $q_\phi(\mathbf{z} \mid \mathbf{a})$ 采用 Transformer 结构，以 CLS token 汇聚序列信息：

\begin{enumerate}[leftmargin=*, label=\textbf{\arabic*.}]
  \item 线性投影：$\mathbf{a} \in \mathbb{R}^{T \times D_a} \xrightarrow{\text{Linear}} \mathbf{X} \in \mathbb{R}^{T \times H_{\text{vae}}}$，$H_{\text{vae}}=256$；
  \item 拼接可学习 CLS token：$[\mathbf{c}_{\text{cls}} \| \mathbf{X}] + \text{PosEmb} \in \mathbb{R}^{(T+1) \times H_{\text{vae}}}$；
  \item $N_{\text{enc}}=3$ 层自注意力 Transformer 块；
  \item 取 CLS 输出经线性层输出 $[\bm{\mu}, \log\bm{\sigma}^2] \in \mathbb{R}^{2d_z}$；
  \item 重参数化采样：
\end{enumerate}
\begin{equation}
  \mathbf{z} = \bm{\mu} + \bm{\varepsilon} \odot \exp\!\left(\tfrac{1}{2}\log\bm{\sigma}^2\right), \quad \bm{\varepsilon} \sim \mathcal{N}(\mathbf{0}, \mathbf{I}_{d_z})
  \label{eq:reparam}
\end{equation}

\subsubsection{ActionVAE：条件解码器}

解码器 $p_\theta(\mathbf{a} \mid \mathbf{z}, \mathbf{F}_{\text{fused}})$ 以 VLM 融合特征为条件，通过跨注意力还原动作序列：

\begin{enumerate}[leftmargin=*, label=\textbf{\arabic*.}]
  \item 线性展开：$\mathbf{z} \xrightarrow{\text{Linear}} \mathbb{R}^{T \times H_{\text{vae}}}$（$T$ 个动作槽）+ 位置编码；
  \item $N_{\text{dec}}=3$ 层解码块，每层依次为：
    \begin{itemize}[leftmargin=*]
      \item 自注意力（动作槽之间）；
      \item \textbf{跨注意力}（动作槽为 Q，$\mathbf{F}_{\text{fused}}$ 为 K/V，注入语言-视觉语义）；
      \item 前馈网络（FFN）；
    \end{itemize}
  \item 线性投影：$\mathbb{R}^{T \times H_{\text{vae}}} \to \hat{\mathbf{a}} \in \mathbb{R}^{T \times D_a}$（输出层零初始化）。
\end{enumerate}

ActionVAE 的训练损失为：
\begin{equation}
  \mathcal{L}_{\text{VAE}} = \underbrace{\|\mathbf{a} - \hat{\mathbf{a}}\|_2^2}_{\mathcal{L}_{\text{recon}}} + \beta \underbrace{D_{\text{KL}}\!\left(q_\phi \,\|\, \mathcal{N}(\mathbf{0},\mathbf{I})\right)}_{\mathcal{L}_{\text{KL}}}
  \label{eq:vae}
\end{equation}
其中 $\beta = 0.001$（$\beta$-VAE 设定），KL 散度正则使隐空间保持连续且近似高斯，便于流匹配采样。

\subsubsection{LatentFlowNet：隐空间速度场}

LatentFlowNet 是一个轻量残差 MLP，在 $d_z=32$ 维隐空间预测流匹配速度场：

\begin{equation}
  v_\theta(\mathbf{z}_t,\, t,\, \mathbf{F}_{\text{fused}},\, \mathbf{h}_K) = \text{MLP}\!\left([\mathbf{z}_t \| \bar{\mathbf{F}}_{\text{fused}} \| \mathbf{h}_K \| \mathbf{e}_t]\right)
  \label{eq:latentfm}
\end{equation}

其中 $\bar{\mathbf{F}}_{\text{fused}} = \text{MeanPool}(\mathbf{F}_{\text{fused}})$ 为 VLM 全局表示，$\mathbf{h}_K$ 为 GRU 编码的本体感知历史，$\mathbf{e}_t$ 为正弦时间编码。网络由 4 层残差块（隐层 512 维，SiLU 激活）组成，输出层零初始化（训练初期速度场近零，稳定收敛）。

隐空间流匹配损失（条件流匹配，Lipman \etal~\cite{flowmatching}）：
\begin{equation}
  \mathcal{L}_{\text{FM}} = \mathbb{E}_{t,\, \mathbf{z}_0 \sim \mathcal{N}(\mathbf{0},\mathbf{I}),\, \mathbf{z}_1 \sim q_\phi} \left\| v_\theta(\mathbf{z}_t, t, \cdot) - (\mathbf{z}_1 - \mathbf{z}_0) \right\|_2^2
  \label{eq:fm}
\end{equation}
其中 $\mathbf{z}_t = (1-t)\mathbf{z}_0 + t\mathbf{z}_1$ 为线性插值路径。

\subsubsection{推理流程}

推理时，LDP 以如下方式生成动作：
\begin{enumerate}[leftmargin=*, label=\textbf{\arabic*.}]
  \item 从标准高斯采样初始噪声：$\mathbf{z}_1 \sim \mathcal{N}(\mathbf{0}, \mathbf{I}_{d_z})$；
  \item Euler 积分 $t: 1 \to 0$（10 步），每步 $\mathbf{z}_{t - \Delta t} = \mathbf{z}_t - \Delta t \cdot v_\theta(\mathbf{z}_t, t, \mathbf{F}_{\text{fused}}, \mathbf{h}_K)$；
  \item 用 ActionVAE 解码器将 $\hat{\mathbf{z}}_0$ 跨注意力到 $\mathbf{F}_{\text{fused}}$，恢复 $\hat{\mathbf{a}} \in \mathbb{R}^{T \times D_a}$；
  \item 反归一化 + gripper 维度二值化（$> 0.5 \to 1$）。
\end{enumerate}

与标准流匹配相比，LDP 的 Euler 积分在 32 维而非 70 维空间进行，每步计算量更小，速度场曲率更低，相同步数下积分误差更小。

\subsection{整体训练目标}

总损失为：
\begin{equation}
  \mathcal{L} = \mathcal{L}_{\text{FM}} + \lambda_{\text{VAE}} \cdot \mathcal{L}_{\text{VAE}}
  \label{eq:total}
\end{equation}
其中 $\lambda_{\text{VAE}} = 1.0$。训练采用 AdamW 优化器，分三组学习率：（1）SmolVLM 骨干：前 1000 步冻结，之后以 $\eta_{\text{base}} \times r_{\text{coef}}$（$r_{\text{coef}} = 0.1$）更新；（2）Action Transformer：$\eta_{\text{base}}$；（3）动作头与 ActionVAE：$\eta_{\text{base}}$。学习率使用线性预热后余弦衰减。

%% =========================================================
\section{实验}
\label{sec:exp}

\subsection{实验设置}

\paragraph{基准}
使用 \textbf{VLABench}~\cite{vlabench2025}，包含 100 类任务：60 类原始任务（Primitive，平均时长~120 步）和 40 类复合任务（Composite，平均时长~500 步）。按 VLABench 标准协议，每类任务从微调集（500 条/类）中采样 100 条轨迹进行评测。

\paragraph{评测指标}
\begin{itemize}[leftmargin=*]
  \item \textbf{成功率（SR）}：episode 级二值成功率；
  \item \textbf{进度分（PS）}：基于子任务完成情况的连续评分；
  \item \textbf{技能召回率（SKR）}：正确识别所需技能的比率；
  \item \textbf{参数召回率（PR）}：正确选取技能参数的比率。
\end{itemize}

\paragraph{实现细节}
所有模型使用 4 块 NVIDIA A100 80GB GPU，fp16 混合精度训练，batch size 32，训练 100k 步（约 \todo{X} 小时）。图像缩放至 $384 \times 384$，动作块长度 $T=10$，ActionVAE 隐变量维度 $d_z=32$，MotionCNN 输出 patch 数 $P=576$。数据增强：颜色抖动、随机裁剪（仅用于静态视角，wrist 视角不增强以保留运动信息）。

\paragraph{基线方法}
\begin{itemize}[leftmargin=*]
  \item \textbf{OpenVLA}~\cite{openvla}：7B 参数，单帧输入，离散 token 动作预测；
  \item \textbf{RDT-1B}~\cite{rdt1b}：1B 参数扩散策略；
  \item \textbf{SimVLA（基线）}~\cite{simvla2026}：0.5B 参数，流匹配，无双流融合；
  \item \textbf{SimVLA + DualStream}：SimVLA + 双流跨注意力融合（无运动引导）；
  \item \textbf{SimVLA-MGA（本文）}：SimVLA + DualStream + MotionCNN + ActionVAE。
\end{itemize}

\subsection{主要结果}

\needtab{表1：VLABench 主结果表。列：方法 | 参数量 | 原始任务 SR↑ | 原始任务 PS↑ | 复合任务 SR↑ | 复合任务 PS↑ | 总体 SR↑ | 总体 PS↑。行：OpenVLA / RDT-1B / SimVLA(baseline) / SimVLA+DualStream / SimVLA-MGA(Ours)。最优用粗体，次优下划线。数据格式：XX.X ± X.X}

\needfig{图3：各方法在 VLABench 6 条评测 Track 上的雷达图（Track 1-6 分别对应：分布内任务/跨类别/常识/语义指令/跨任务/未见纹理）。}

如表~\ref{tab:main} 所示（\todo{填入实验数据后取消占位符}），SimVLA-MGA 在总体成功率上超越所有基线，在复合任务上的提升尤为显著。相比 SimVLA + DualStream 基线：

\begin{itemize}[leftmargin=*]
  \item \textbf{运动引导注意力}（MGCA）在原始任务上平均提升 \todo{X.X\%} SR，在复合任务上提升 \todo{X.X\%} SR，说明腕部运动信息对长时序阶段切换有重要作用；
  \item \textbf{ActionVAE} 进一步提升复合任务 PS \todo{X.X}，表明隐空间流匹配降低了多阶段动作序列的预测方差；
  \item 两者\textbf{协同增益}超过各自单独贡献之和，验证了运动感知与隐式策略压缩的互补性。
\end{itemize}

\subsection{消融研究}

\needtab{表2：消融实验。逐步添加各组件：(a) SimVLA baseline；(b) + DualStream；(c) + DualStream + MGCA；(d) + DualStream + ActionVAE；(e) 完整 SimVLA-MGA。指标：原始 SR / 复合 SR / 总体 SR / 参数增量(M)}

表~\ref{tab:ablation} 的消融结果揭示以下规律：

\textbf{双流融合的基础作用。} 单独添加 DualStream 即带来约 \todo{X.X\%} SR 提升，确认了将 wrist 特征与静态视角特征分流融合的有效性。

\textbf{MGCA 的增益来源于运动区域聚焦。} 我们进一步设计了对照实验：将 MotionCNN 替换为随机噪声激活图（random map）。如表~\ref{tab:ablation} 所示，随机图不带来额外增益（甚至轻微下降），而真实帧差分驱动的 MGCA 稳定提升，说明性能增益确实来自运动语义而非注意力扰动本身。

\textbf{ActionVAE 对复合任务的针对性提升。} 在原始任务（平均时长~120 步）上 ActionVAE 的增益有限（约 \todo{X.X\%}），而在复合任务（平均时长~500 步）上增益显著（约 \todo{X.X\%}），符合"隐空间压缩降低长时序累积方差"的理论预期。

\textbf{可学习偏置缩放因子 $\alpha$ 的重要性。} 若将 $\alpha$ 固定为 1（不可学习），性能相比零初始化可学习版本下降约 \todo{X.X\%}，说明让模型自主决定运动信号强度的必要性。

\subsection{定性分析}

\needfig{图4：典型成功/失败案例对比。选取 VLABench 中 3 个复合任务（如：抓取特定物体→移动至目标位置→精确放置）。每行：任务描述 | 关键帧序列（4帧）| SimVLA+DualStream 结果 | SimVLA-MGA 结果 | 对应运动激活图热力图。标注成功（✓）/失败（✗）及失败原因。}

\needfig{图5：运动偏置缩放因子 $\alpha$ 在训练过程中的变化曲线。展示 $\alpha$ 从 0 逐步增大并收敛的动态过程，以及与 training loss 曲线的对应关系。}

\paragraph{运动激活图分析}
如图~\ref{fig:motion_vis} 所示，在抓取阶段，MotionCNN 准确激活了末端执行器与目标物体的接触区域；在移动阶段，激活区域随末端执行器轨迹移动；在放置阶段，激活集中于物体即将接触目标位置的区域。这一模式表明 MotionCNN 确实学到了有意义的运动语义，而非简单响应相机噪声。

\paragraph{失败模式分析}
SimVLA-MGA 的主要失败案例集中在：（1）强光照变化导致帧差分产生误导性激活（相机伪运动）；（2）任务需要精确的语言理解而视觉运动信息不足（如"拿起标有'X'的物品"）。这两类问题为未来工作指出了方向。

\subsection{效率分析}

\needtab{表3：效率对比。列：方法 | 参数量(B) | 训练显存(GB, batch=32) | 推理时延(ms/step) | 训练时长(GPU·hours)}

SimVLA-MGA 相比 SimVLA 基线仅新增约 3M 参数（< 0.7\%），推理时延增加约 \todo{X} ms/步（MotionCNN 的前向开销），在效率与性能之间取得良好平衡。

%% =========================================================
\section{结论}
\label{sec:conclusion}

本文提出 SimVLA-MGA，在轻量 VLA 基线之上以极小参数代价引入两项互补创新：运动引导跨视角注意力机制（MGCA）和 ActionVAE 隐式流匹配策略。前者通过帧差分驱动的 MotionCNN 将运动语义注入跨注意力，使静态视角特征自动聚焦于腕部运动区域；后者将动作序列压缩至低维隐空间再执行流匹配，有效降低长时序策略的预测方差。在 VLABench 100 类任务上的全面评测验证了两者的有效性及其协同增益，在复合长时序任务上的显著提升说明运动感知对多阶段机器人操作的重要价值。

\paragraph{局限性与未来工作}
当前方法对相机伪运动（光照变化、轻微抖动）存在一定敏感性；未来可探索在差分图上引入光流补偿或利用 IMU 信号加以校正。此外，将 MGCA 扩展至更多视角（如双腕部相机设置）也是值得探索的方向。

%% =========================================================
\section*{致谢}
\todo{填写致谢内容（基金项目编号、计算资源来源等）。}

%% =========================================================
%% 参考文献
\bibliographystyle{plain}
% 若使用 BibTeX，取消下一行注释并创建 refs.bib
% \bibliography{refs}

% 如未使用 BibTeX，手动列出参考文献如下：
\begin{thebibliography}{99}

\bibitem{simvla2026}
Y.~Luo, W.~Chen, T.~Liang, B.~Wang, and Z.~Li.
\newblock {SimVLA}: A Simple VLA Baseline for Robotic Manipulation.
\newblock \textit{arXiv preprint arXiv:2602.18224}, 2026.

\bibitem{vlabench2025}
S.~Zhang, Z.~Xu, P.~Liu, X.~Yu, Y.~Li, Q.~Gao, Z.~Fei, Z.~Yin, Z.~Wu, Y.-G.~Jiang, and X.~Qiu.
\newblock {VLABench}: A Large-Scale Benchmark for Language-Conditioned Robotics Manipulation with Long-Horizon Reasoning Tasks.
\newblock In \textit{ICCV}, 2025.

\bibitem{openvla}
M.~J. Kim, K.~Pertsch, S.~Karamcheti, T.~Xiao, A.~Balakrishna, S.~Nair, et al.
\newblock {OpenVLA}: An Open-Source Vision-Language-Action Model.
\newblock \textit{arXiv preprint arXiv:2406.09246}, 2024.

\bibitem{rt2}
A.~Brohan, N.~Brown, J.~Carbajal, Y.~Chebotar, X.~Chen, K.~Choromanski, et al.
\newblock {RT-2}: Vision-Language-Action Models Transfer Web Knowledge to Robotic Control.
\newblock In \textit{CoRL}, 2023.

\bibitem{pi0}
K.~Black, N.~Brown, D.~Driess, A.~Esmail, M.~Equi, C.~Finn, et al.
\newblock $\pi_0$: A Vision-Language-Action Flow Model for General Robot Control.
\newblock \textit{arXiv preprint arXiv:2410.24164}, 2024.

\bibitem{roboflamingo}
J.~Li, D.~Zhang, H.~Liu, S.~Liu, B.~Liu, H.~Dong, Y.~Qiao, and C.~Gao.
\newblock {RoboFlamingo}: Vision-Language Models as Prefix Learners for Robot Manipulation.
\newblock \textit{arXiv preprint arXiv:2311.01378}, 2023.

\bibitem{diffpolicy}
C.~Chi, S.~Feng, Y.~Du, Z.~Xu, E.~Cousineau, B.~Burchfiel, and S.~Song.
\newblock Diffusion Policy: Visuomotor Policy Learning via Action Diffusion.
\newblock In \textit{RSS}, 2023.

\bibitem{flowmatching}
Y.~Lipman, R.~T.~Q. Chen, H.~Ben-Hamu, M.~Nickel, and M.~Le.
\newblock Flow Matching for Generative Modeling.
\newblock In \textit{ICLR}, 2023.

\bibitem{rold}
Y.~Zhou, J.~Chen, Y.~Hu, Y.~Zhang, J.~Gao, and M.~Liu.
\newblock {RoLD}: Robot Learning from Demonstrations via Latent Diffusion.
\newblock \textit{arXiv preprint arXiv:2403.07312}, 2024.

\bibitem{deltacnn}
L.~Parger, C.~Wang, T.~Reinert, and M.~Wimmer.
\newblock {DeltaCNN}: End-to-End CNN Inference of Sparse Frame Differences in Videos.
\newblock In \textit{CVPR}, 2022.

\bibitem{motiondeltacnn}
L.~Parger, C.~Wang, T.~Reinert, and M.~Wimmer.
\newblock {MotionDeltaCNN}: Sparse CNN Inference of Frame Differences in Moving-Camera Videos with Continuous Motion Fields.
\newblock In \textit{ICCV}, 2023.

\bibitem{rdt1b}
S.~Liu, H.~Chen, B.~Bai, W.~Feng, H.~Ouyang, L.~Wang, W.~Chen, K.~Cui, and S.~Ji.
\newblock {RDT-1B}: A Diffusion Foundation Model for Bimanual Manipulation.
\newblock \textit{arXiv preprint arXiv:2410.07864}, 2024.

\end{thebibliography}

\end{document}
